%% paper.tex — The Guarded Syntax
%% Gate: scripts/paper-check.sh (label + discipline checks)
%%       scripts/sync-check.sh  (Dcare.* names exist in the mechanization)
%%       latexmk -pdf paper.tex  (clean build)
\documentclass{eptcs}
\usepackage{guarded-syntax}

%% Contact email deliberately absent: the ORCID iD is verified against the
%% program's Zenodo deposits, and no email is on record, so none is invented.
\title{The Guarded Syntax: A String-Diagram Notation\\ for Linear Resource Calculi}
\author{
  Jaye Timothy Marshall%
  \institute{Independent researcher\\
             {\footnotesize ORCID 0009-0002-5209-8161}}
}
\def\authorrunning{J.\,T. Marshall}
%% EPTCS defines no \keywords command; none is invented.
\begin{document}
\maketitle
\begin{abstract}
We introduce the \emph{Guarded Syntax}, a written notation for the Care
Calculus, a multi-sorted ordered linear logic with coalgebraic
generators and guarded recursion, mechanized in Lean~4. The notation has
five operators: the knower $\gknower$, the one Cartesian source; the
linear wire $\glto$, which consumes the left and produces the right
exactly once; the guard $\guard$, one interval of delay; the feed
$\gfeed$, which binds a delayed output back into its own definition (the
$+1$); and the sink $\gsink$, the archive where content is recorded and
never read (the $-1$). Agent definitions are one-line expressions of
this grammar, and they have a planar string-diagram reading --- wires in
two gauges, guard beads, feedback envelopes --- that we develop
alongside: the standard diagram languages build in a symmetry that
ordered resource calculi do not have, so the reading is planar, and the
one definable exchange is an explicit swap-box. The calculus's two
linearity laws appear twice: as expressions that fail to type (a
cloner, an eraser --- the archive is legal, and the difference is
load-bearing), and as theorem statements with machine-checked fragment
strengths. We state the correspondence precisely: well-formed
expressions translate to derivations and back (claims with proof
sketches, anchored at the term level to machine-checked theorems); the
linear cut is the splicing of a pending wire into the hole it closes,
resolving the ``mutual promise'' of two abstractions each waiting on
the other as an unreduced composite; and feedback is legal exactly when
guarded --- the diagrammatic form of the side condition that makes
self-application cost a context. Well-formedness is compositional:
there is no global correctness criterion, and the two linearity laws
are not forbidden moves to be checked after writing but shapes the
grammar simply cannot generate --- a linear wire has no junction, and a
content wire has no bare end. Every diagram in this paper is drawn from
a named, machine-checked derivation; the notation-level correspondence
is stated, not mechanized.
\end{abstract}

\section{Introduction}\label{sec:intro}

The Care Calculus is a multi-sorted ordered linear logic with three
features that ordinary linear logic does not carry: a coalgebraic
generator at a distinguished non-linear sort, a delay modality $\guard$
with guarded recursion, and an archive type $\tT$ into which content can
be recorded but from which nothing can be read. It was developed to study
when protection for a party can be \emph{derived from the structure} of an
optimizing system rather than trained in, and it is fully mechanized: a
Lean~4 development (\texttt{Dcare}, several hundred machine-checked
theorem declarations) proves its linearity discipline, its guardedness
properties, and the non-derivability facts that make the calculus a
resource grammar rather than a suggestion \cite{dcarecompanion}.

The mechanization is precise and, in textual form, hard to read. The
resource discipline of a split-context judgment --- \emph{who} is
consumed, \emph{what} is delayed, \emph{which} wire loops --- is scattered
across context equations $\Gamma\,{+}\,\Delta = E$ that a proof assistant
checks happily and a human tracks poorly. This paper introduces a
notation, the \emph{Guarded Syntax}, in which those three questions become
inspection: consumption is a wire's gauge, delay is a bead, feedback is an
envelope.

Existing graphical languages do not fit, for one structural reason each:

\begin{itemize}
\item \textbf{Proof nets} \cite{girard1987,danos1989} are the geometry of
multiplicative linear logic, but their well-formedness is a \emph{global}
correctness criterion on a graph, their wires are undirected, and they
have no native place for a delay or a loop.
\item \textbf{Symmetric monoidal string diagrams} \cite{joyal1991} make
wire-crossing free: they present categories with exchange. The Care
Calculus is \emph{ordered} --- its contexts split in sequence and it has
no structural exchange --- so the symmetric notation would identify
derivations the judgment distinguishes.
\item \textbf{Signal-flow graphs as props} \cite{bsz2014} are the closest
kin in spirit --- directed wires, boxes, feedback --- but their wires
fan out freely, which is exactly what a linear resource may not do.
\end{itemize}

The Guarded Syntax is therefore \emph{planar} \cite{polakow2017}: wires
never cross for free, and the one definable exchange of the calculus (a
$\mathit{pair}$--$\mathit{lete}$ round trip) is drawn as an explicit
swap-box rather than as isotopy. Its contributions:

\begin{itemize}
\item a \textbf{written grammar} of five operators --- $\gknower$,
$\glto$, $\guard$, $\gfeed$, $\gsink$ --- whose typed expressions are
one-line agent definitions, with the two linearity laws as type errors
(\S\ref{sec:core});
\item its \textbf{planar string-diagram reading} --- two wire gauges,
guard beads, feedback envelopes --- in which the calculus's resource
discipline is visible on the page (\S\ref{sec:core:reading});
\item the \textbf{Care signature} as labelled boxes, with the laws'
theorem statements and fragment strengths (\S\ref{sec:signature},
\S\ref{sec:laws});
\item \textbf{the correspondence stated precisely} --- five claims with
proof sketches, each anchored at the term level to a named machine-checked
theorem (\S\ref{sec:correspondence});
\item two \textbf{worked readings}: the epistemic cycle, written as one
line and drawn from a machine-checked witness term; and the cut, drawn
from a machine-checked substitution principle (\S\ref{sec:worked}).
\end{itemize}

What this paper is not: it is not new type theory --- every term-level
theorem cited here is proven in \texttt{Dcare} and cited by name, and the
two unrestricted classical statements of the calculus's linearity laws
remain open conjectures there, proven in fragment strengths; and it is not
a mechanization of the diagram translation --- that correspondence is
stated as claims, not checked by a machine (\S\ref{sec:correspondence}).
Interpretive remarks are confined to \S\ref{sec:remarks} and are labeled.

\paragraph{Artifact availability.}
The Lean~4 development \texttt{Dcare} and the complete \LaTeX{} source of
this paper, including the diagram macros of \texttt{guarded-syntax.sty},
are available under GNU AGPL-3.0.

\section{The Guarded Syntax as a written grammar}\label{sec:core}

The Guarded Syntax is, first, a notation for \emph{writing}: five
operators that compose into typed expressions, in which a whole agent
definition fits on one line. Its two-dimensional reading --- glyphs laid
out as wires on a page --- is derived afterwards
(\S\ref{sec:core:reading}): the symbols are the syntax, the picture is a
layout of the syntax.

\subsection{The five operators}\label{sec:core:ops}

\begin{center}
\begin{tabular}{clp{60mm}}
\textbf{glyph} & \textbf{name} & \textbf{role in the grammar}\\
\hline
$\gknower$ & the knower & the Cartesian source: the one sort that may be
copied and discarded; every content token descends from one\\
$\glto$ & the linear wire & the type $\tau \glto \sigma$: consumes the
left, produces the right, exactly once\\
$\guard$ & the guard & the type $\guard\tau$ and the prefix on recursive
definitions: one interval of delay\\
$\gfeed$ & the feed & the feedback operator $e \gfeed x$, binding a
delayed output back into its own definition --- the $+1$\\
$\gsink$ & the sink & the archive constant and its type: content is
recorded there and never read --- the $-1$ debt\\
\end{tabular}
\end{center}

The judgment is $\Gamma \vdash e : \tau$ --- ``$e$ realizes $\tau$ from
the resources $\Gamma$'' --- with types and expressions:

\begin{center}
\fbox{\begin{minipage}{0.94\linewidth}\centering\vspace{3pt}
$\begin{array}[t]{rcl}
\tau & ::= & \tC \mid \tK \mid \tA \mid \tV \mid \tS \mid \tT \mid \tUnit
   \mid \tau \glto \tau \mid \guard\tau \mid \tau \ten \tau\\[3pt]
e & ::= & \gknower \mid x \mid \lambda x{:}\tau.\,e \mid e\,e
   \mid e \ten e \mid \mathit{lete}\;e\;e \mid () \mid \guard e \mid e \gfeed x\\
  & \mid & \mathit{cons}(e,e) \mid \mathit{nil} \mid \mathit{snoc}(e,e)
   \mid \mathit{gen} \mid \mathit{red} \mid \mathit{obj} \mid \mathit{terminate}\\
\end{array}$\vspace{3pt}
\end{minipage}}
\end{center}

Every production re-skins a rule of the mechanized judgment: $\glto$ is
its arrow at linear sorts, $\guard$ its delay, $e \gfeed x$ its guarded
fixpoint (one lap of the feed consumes exactly one guard --- the $+1$),
$\gknower$ marks a hypothesis of the knower sort, and $\gsink$
abbreviates the archive constructors $\mathit{nil}$ and $\mathit{snoc}$.
The notation adds no rule and removes none; that is the content of the
correspondence stated in \S\ref{sec:correspondence}.

\subsection{The two laws, as type errors}\label{sec:core:laws}

The brief's two forbidden shapes are literally expressions that fail to
type --- and the third line below is the surprise the mechanization
forces:

\begin{center}
\fbox{\begin{minipage}{0.94\linewidth}\vspace{3pt}
{\small $\mathit{clone} : \tC \glto \tC \ten \tC$}
\hfill \lawerr\ no rule copies --- contraction exists only at
$\gknower$\vspace{3pt}\\
{\small $\mathit{erase} : \tC \glto \tUnit$}
\hfill \lawerr\ a consumed token must be carried, reduced, or archived
--- not dropped\vspace{3pt}\\
{\small $\mathit{archive} : \tC \glto \gsink$}
\hfill \checkmark\ \emph{legal}: the $-1$ debt, booked forever\vspace{3pt}
\end{minipage}}
\end{center}

Archiving into $\gsink$ is derivable (\dcn{archival_witness}
\meas{machine-checked}), so the sink is not an error site but the
calculus's most loaded legal shape: the count certifies the token was
spent, nothing certifies it was heard. The forbidden vanishing shape is
the bare end ($\glto\,\tUnit$), not the sink. Section~\ref{sec:laws}
carries the full statement table with fragment strengths.

\subsection{The pictorial reading}\label{sec:core:reading}

% fig-core.tex — the core graphical language (gallery).
\begin{figure}[t]
\centering
\begin{tikzpicture}[node distance=4mm]
  % (1) linear wire
  \node[lbl] (l1) at (0,0) {linear wire $\glto$:};
  \draw[wire] (5.6,0) -- (8.0,0);
  \node at (6.8,0.22) {$\glto$};
  % (2) non-linear rail
  \node[lbl] (l2) at (0,-0.75) {non-linear rail (the thick rail):};
  \draw[thickwire] (5.6,-0.75) -- (8.0,-0.75);
  % (3) box
  \node[lbl] (l3) at (0,-1.6) {box (signature constant):};
  \node[box] (b3) at (6.8,-1.6) {$\mathit{red}$};
  \draw[wire] (5.6,-1.6) -- (b3.west);
  \draw[wire] (b3.east) -- (8.0,-1.6);
  % (4) guard bead
  \node[lbl] (l4) at (0,-2.45) {guard bead ($\guard$) on a wire:};
  \draw[wire] (5.6,-2.45) -- (6.7,-2.45);
  \node[bead] (be4) at (6.9,-2.45) {};
  \draw[wire] (7.1,-2.45) -- (8.0,-2.45);
  % (5) guarded feedback envelope
  \node[lbl] (l5) at (0,-3.7) {guarded feedback ($\gfeed$):};
  \node[box] (b5) at (7.0,-3.45) {};
  \draw[wire] (6.1,-3.45) -- (b5.west);
  \draw[wire] (b5.east) -- (8.4,-3.45);
  \node[bead] (be5) at (7.35,-4.35) {};
  \draw[wire] (7.75,-3.72) to[out=-90,in=0] (7.5,-4.35) -- (be5);
  \draw[wire] (be5) -- (6.7,-4.35) to[out=180,in=-90] (6.4,-3.72);
  \node[env,fit=(b5)(be5),inner sep=5pt] (e5) {};
  % (6) sink
  \node[lbl] (l6) at (0,-5.1) {archive sink ($\gsink$):};
  \draw[wire] (5.6,-5.1) -- (7.0,-5.1);
  \node[sink] at (7.0,-5.18) {};
  \draw[thick] (6.9,-5.32) -- (7.1,-5.32);
  \draw[thick] (6.94,-5.42) -- (7.06,-5.42);
  % (7) swap-box
  \node[lbl] (l7) at (0,-5.95) {swap-box ($\gswap$):};
  \node[box,minimum width=5mm] (sw) at (6.8,-5.95) {$\gswap$};
  \draw[wire] (5.6,-5.7) -- (sw.north west);
  \draw[wire] (5.6,-6.2) -- (sw.south west);
  \draw[wire] (sw.north east) -- (8.0,-6.2);
  \draw[wire] (sw.south east) -- (8.0,-5.7);
\end{tikzpicture}
\caption{The core graphical language. Wires are directed and carry a gauge:
thin wires are linear (used exactly once), thick rails are non-linear (may
branch and may be discarded). Boxes are the signature's constants. The guard
$\guard$ is a decorator \emph{on} a wire, not a box; a feedback envelope is a
dashed region whose internal loops must pass through at least one guard. The
sink $\gsink$ is the archive terminal. The swap-box is the only crossing
artifact, and it is a defined composite, not free isotopy
(\S\ref{sec:core:swap}).}
\label{fig:core}
\end{figure}

Written expressions have a two-dimensional reading in which each glyph
becomes a drawing element (Fig.~\ref{fig:core}); the rest of this paper
uses both forms freely. A $\glto$-type is drawn as a directed thin wire;
a $\gknower$-type as a thick rail --- the gauges record use-once versus
copy-at-will, the discipline of the counting lemma
(\dcn{single_linear_use} \meas{machine-checked}). $\guard$ decorates the
wire it delays: a bead. $e \gfeed x$ is a dashed envelope around $e$
whose loop carries the delayed binder back through a bead. $\gsink$ is
the terminal. Boxes carry the named constants of the signature
(\S\ref{sec:signature}); juxtaposition is ordered tensor --- \emph{which}
side a box sits on is information, not layout.

\subsection{The guard bead}\label{sec:core:bead}

The $\guard$ \emph{guard} is a decorator on a wire, never a box. One bead
marks one interval of delay: a wire $\guard\tau$ carries $\tau$ one step
later than the present. Beads are the notation's clock; they enter the
calculus twice, once at the stream constructor --- a guarded stream has a
head now and a $\guard$-delayed tail --- and once at recursion
(\S\ref{sec:core:feedback}). At the term level, the one-step unfolding of
a guarded fixpoint consumes exactly one bead and preserves typing (the
theorem \dcn{fix_unfold_closed} \meas{machine-checked}). The
bead is the diagrammatic kin of the later modality of guarded recursion
\cite{nakano2000,birkedal2012}.

\subsection{Guarded feedback}\label{sec:core:feedback}

\emph{Feedback} $\gfeed$ is drawn as a dashed \emph{envelope}: a region of
the diagram containing a loop from an output wire back to an input port.
Two laws govern it, and both are load-bearing:

\begin{enumerate}
\item \textbf{The guard law.} A feedback loop must pass through at least
one guard bead. A loop with no bead is not a legal envelope --- it is an
unreduced composite (\S\ref{sec:worked:cut} shows one). This is the
diagrammatic form of guardedness: the bead is what makes the recursive
binding well-founded, one interval per lap.
\item \textbf{The permeability law.} Only thick rails may cross the
envelope boundary. The recursive binding closes over the non-linear
context and nothing else: at the term level this is exactly the side
condition of the guarded fixpoint rule, which requires the context to be
all non-linear --- the price of self-application made syntactic (the
witness context fact \dcn{knower_ctx_nonlinear}, and the typed
counterexample \dcn{selfapp_price}, are
\meas{machine-checked}).
\end{enumerate}

A comment on what a loop \emph{is} here. The feedback curve departs from
an output wire, and to an eye trained on Cartesian circuits that looks
like a tap. It is not one: a tap is a junction, and thin wires have no
junctions. The loop is a \emph{recursive binding} --- the fixed-point
convention of synchronous dataflow \cite{lustre1991}, read through a trace
\cite{jsv1996} that exists only because the guard makes the equation
causal. What the notation draws as a curve closing a loop, the judgment
writes as a fixpoint binding a delayed variable.

\subsection{The swap-box}\label{sec:core:swap}

Wires do not cross for free. The one crossing artifact is the
\emph{swap-box} $\gswap$ (Fig.~\ref{fig:swap}): a two-in, two-out box that
rebinds two adjacent wires in the opposite order. It is not structure ---
the calculus has no exchange rule --- it is a \emph{definable} composite,
namely a $\mathit{pair}$ followed by a $\mathit{lete}$, exhibited inside
the machine-checked witness \dcn{knower_sustains_stream}
\meas{machine-checked}. Planar isotopy is free; commutation costs a box.
This is the single sharpest difference between the Guarded Syntax and the
sym-monoidal habit: what braiding gives for free, the ordered world must
build, and the notation shows the price on the page.
% fig-swap.tex — the definable swap-box.
\begin{figure}[t]
\centering
\begin{tikzpicture}
  % left: the composite (pair then lete) — the derived internal exchange
  \node[box] (pr) at (0,0) {$\mathit{pair}$};
  \node[box] (lt) at (2.6,0) {$\mathit{lete}$};
  \draw[wire] (-1.6,0.45) -- (pr.north west);
  \node[lbl] at (-1.85,0.75) {$\tS_1$};
  \draw[wire] (-1.6,-0.45) -- (pr.south west);
  \node[lbl] at (-1.85,-0.75) {$\tK_2$};
  \draw[wire] (pr.east) -- (lt.west);
  \node[lbl] at (1.3,0.22) {$\tS\ten\tK$};
  \draw[wire] (lt.east) -- (4.0,0.45);
  \node[lbl] at (4.15,0.75) {$\tK$};
  \draw[wire] (lt.east) -- (4.0,-0.45);
  \node[lbl] at (4.15,-0.75) {$\tS$};
  % equals
  \node[font=\Large] at (5.2,0) {$=$};
  % right: the swap-box abbreviation
  \node[box,minimum width=5.5mm] (sw) at (6.6,0) {$\gswap$};
  \draw[wire] (5.8,0.45) -- (sw.north west);
  \node[lbl] at (5.55,0.75) {$\tS$};
  \draw[wire] (5.8,-0.45) -- (sw.south west);
  \node[lbl] at (5.55,-0.75) {$\tK$};
  \draw[wire] (sw.north east) -- (7.6,-0.45);
  \node[lbl] at (7.75,-0.75) {$\tS$};
  \draw[wire] (sw.south east) -- (7.6,0.45);
  \node[lbl] at (7.75,0.75) {$\tK$};
\end{tikzpicture}
\caption{The swap-box: the calculus has no structural exchange, but a
$\mathit{pair}$--$\mathit{lete}$ round trip rebinds two adjacent hypotheses
in the opposite order — a \emph{definable} exchange, exhibited as the
internal flip inside the machine-checked witness
\texttt{Dcare.knower\_sustains\_stream}. The Guarded Syntax draws it as the
explicit two-port box $\gswap$, never as free wire-crossing: planar isotopy
is free, commutation is not.}
\label{fig:swap}
\end{figure}


\subsection{Sinks}

The $\gsink$ \emph{sink} is a terminal for a wire: what enters is
recorded, not destroyed and not read. In the Care Calculus it draws the
archive type $\tT$, which has no elimination rule --- the past is inert;
no term reduces a trace to knowledge (the theorem
\dcn{past_cannot_teach} \meas{machine-checked}). The symbol is
$\bot$ by lineage, not by meaning: it is not falsity, and reading it as
``content goes to die'' is exactly right --- dying is terminal, and legal
(\S\ref{sec:laws}).

\subsection{Well-formedness is compositional}

A diagram is well-formed by construction: the grammar generates it from
wires, boxes, beads, envelopes and swap-boxes, and every local
construction is type-correct. There is no global correctness criterion to
check afterwards, no switching condition, no acyclicity test beyond the
local guard law on each envelope. This is a deliberate divergence from
proof nets \cite{danos1989}, whose correctness is a property of a graph
audited after drawing. The practical consequence: an ill-shaped diagram
--- a cloned wire, a vanished resource --- cannot be drawn at all, so it
cannot be mistaken for a proof.

\section{The Care signature}\label{sec:signature}

The Care Calculus has seven sorts --- content $\tC$, knower $\tK$,
knowledge $\tA$, value $\tV$, stream $\tS$, trace $\tT$, unit $\tUnit$ ---
and the delay $\guard\tau$ and tensor $\tau\ten\tau$ on all of them.
Exactly one sort is non-linear: $\tK$, the \emph{knower}. Everything else
is linear. The signature (Fig.~\ref{fig:signature}) has seven constants:
% fig-signature.tex — the Care signature as boxes.
\begin{figure}[t]
\centering
\begin{tikzpicture}
  % --- gen : K -> C (x) K ---
  \node[box] (gen) at (0,0) {$\mathit{gen}$};
  \node[circle,draw,thick,fill=black!12,inner sep=0pt,minimum size=3.2mm]
       (kn) at (-1.7,0) {$\cdot$};
  \node[lbl] at (kn) [below=0mm] {$\gknower\;\tK$};
  \draw[thickwire] (kn) -- (gen.west);
  \draw[wire] (gen.east) ++(0,0) -- ++(0.5,0) -- ++(0,0.5);
  \node[lbl] at (1.15,0.62) {$\tC$};
  \draw[thickwire] (gen.east) ++(0,0) -- ++(0.5,0) -- ++(0,-0.5);
  \node[lbl] at (1.15,-0.62) {$\tK$};
  % --- red : S -> A ---
  \node[box] (red) at (5.2,0) {$\mathit{red}$};
  \draw[wire] (4.2,0) -- (red.west);
  \node[lbl] at (4.05,0.22) {$\tS$};
  \draw[wire] (red.east) -- (6.2,0);
  \node[lbl] at (6.35,0.22) {$\tA$};
  % --- obj : S -> V ---
  \node[box] (obj) at (5.2,-1.5) {$\mathit{obj}$};
  \draw[wire] (4.2,-1.5) -- (obj.west);
  \node[lbl] at (4.05,-1.28) {$\tS$};
  \draw[wire] (obj.east) -- (6.2,-1.5);
  \node[lbl] at (6.35,-1.28) {$\tV$};
  % --- cons : C, ▷S -> S ---
  \node[box] (cons) at (0,-3.4) {$\mathit{cons}$};
  \draw[wire] (-1.4,-3.0) -- (cons.north west);
  \node[lbl] at (-1.55,-2.78) {$\tC$};
  \draw[wire] (-1.4,-3.8) -- (cons.south west);
  \node[bead] (cb) at (-1.95,-3.8) {};
  \draw[wire] (cb) -- (-1.4,-3.8);
  \node[lbl] at (-2.1,-3.55) {$\guard\tS$};
  \draw[wire] (cons.east) -- (1.2,-3.4);
  \node[lbl] at (1.35,-3.18) {$\tS$};
  % --- nil : 1 -> T ---
  \node[box] (nil) at (5.2,-3.4) {$\mathit{nil}$};
  \draw[wire] (nil.east) -- (6.2,-3.4);
  \node[lbl] at (6.35,-3.18) {$\tT$};
  \node[lbl] at (4.6,-3.15) {$\tUnit$};
  \draw[wire] (4.2,-3.4) -- (nil.west);
  % --- snoc : T, C -> T ---
  \node[box] (snoc) at (0,-5.6) {$\mathit{snoc}$};
  \draw[wire] (-1.4,-5.25) -- (snoc.north west);
  \node[lbl] at (-1.55,-5.03) {$\tT$};
  \draw[wire] (-1.4,-5.95) -- (snoc.south west);
  \node[lbl] at (-1.55,-6.2) {$\tC$};
  \draw[wire] (snoc.east) -- (1.2,-5.6);
  \node[lbl] at (1.35,-5.38) {$\tT$};
  % --- terminate : K -> T ---
  \node[box] (term) at (5.2,-5.6) {$\mathit{terminate}$};
  \draw[thickwire] (4.2,-5.6) -- (term.west);
  \node[lbl] at (4.05,-5.38) {$\tK$};
  \draw[wire] (term.east) -- (6.2,-5.6);
  \node[lbl] at (6.35,-5.38) {$\tT$};
\end{tikzpicture}
\caption{The Care signature. Each constant of the calculus is a box with
typed ports; thin wires are linear sorts, the thick rail is the knower
$\gknower$~$\tK$. The generator $\mathit{gen}$ is the coalgebra at the
knower node; $\mathit{cons}$ demands a delayed tail ($\guard\tS$ — the bead
on its second input); $\mathit{nil}$ and $\mathit{snoc}$ build and extend
the archive $\tT$; $\mathit{terminate}$ archives a knower in one step.}
\label{fig:signature}
\end{figure}


\begin{itemize}
\item $\mathit{gen} : \tK \glto \tC \ten \tK$ --- the coalgebra at the
knower: firing it on a knower yields fresh content \emph{and the knower
back}. The knower node $\gknower$ marks this: the source whose rail
persists through each firing.
\item $\mathit{red} : \tS \glto \tA$ and $\mathit{obj} : \tS \glto \tV$
--- reduction and valuation, the two consumers of the live stream. They
consume the whole infinite stream, not a prefix.
\item $\mathit{cons} : \tC \ten \guard\tS \glto \tS$ --- the guarded
stream constructor: head now, tail one interval later. The bead on its
second input is the clock tick.
\item $\mathit{nil} : \tUnit \glto \tT$ and $\mathit{snoc} : \tT \ten \tC
\glto \tT$ --- the archive constructors: a trace is built from nothing and
extended only at its end.
\item $\mathit{terminate} : \tK \glto \tT$ --- archiving a knower in one
step; the terminal event.
\end{itemize}

Reading the signature as a diagram is already informative. Two of the
seven constants produce content ($\mathit{gen}$, from a knower; nothing
else does), one consumes it into the archive ($\mathit{snoc}$), and the
stream side is the only place where delay is constructive. The shape the
signature cannot express --- a constant $\tC \glto \tC \ten \tC$, or
$\tC \glto \tUnit$ --- is the subject of the next section.

\section{The two laws as ungeneratable shapes}\label{sec:laws}

The calculus's two classical linearity statements are non-derivability
facts: contraction for content, $[\tC] \nvdash \tC \ten \tC$, and
vanishing, $[\tC] \nvdash \tUnit$. In the Guarded Syntax they are not
rules to check; they are absences in the grammar
(Fig.~\ref{fig:laws}):
% fig-laws.tex — the two laws as non-generatable shapes (negative pedagogy).
\begin{figure}[t]
\centering
\begin{tikzpicture}
  % ---- panel (a): no cloning ----
  \node[lbl,align=left] at (1.5,1.3)
    {(a) a thin wire has\\ \emph{no junction}};
  \draw[wire] (0.0,0.6) -- (1.5,0.6);
  \draw[wire] (1.5,0.6) -- (2.7,1.2);
  \draw[wire] (1.5,0.6) -- (2.7,0.0);
  \node[font=\Large,text=red!80!black] at (1.5,0.62) {$\times$};
  \node[lbl] at (3.35,1.25) {$\tC$};
  \node[lbl] at (3.35,0.05) {$\tC$};
  \node[lbl] at (0.0,0.85) {$\tC$};
  % the legal contrast: the knower rail branches
  \node[circle,draw,thick,fill=black!12,inner sep=0pt,minimum size=3mm]
       (kn) at (0.75,-1.3) {};
  \node[lbl] at (0.75,-1.62) {$\gknower\;\tK$};
  \draw[thickwire] (kn.east) -- (1.9,-0.85);
  \draw[thickwire] (kn.east) -- (1.9,-1.75);
  \node[lbl] at (2.15,-0.8) {$\tK$};
  \node[lbl] at (2.15,-1.85) {$\tK$};
  \node[lbl,text=black!45] at (1.5,-2.35) {legal: rails branch};
  % ---- panel (b): no vanishing ----
  \node[lbl,align=left] at (7.6,1.3)
    {(b) a content wire has\\ \emph{no bare end}};
  \draw[wire] (6.0,0.6) -- (7.6,0.6);
  \node[font=\Large,text=red!80!black] at (7.95,0.6) {$\times$};
  \node[lbl] at (6.0,0.85) {$\tC$};
  \node[lbl,text=black!45] at (7.6,0.15) {the wire just stops};
  % legal: archived
  \draw[wire] (6.0,-0.7) -- (7.3,-0.7);
  \node[sink] at (7.3,-0.78) {};
  \draw[thick] (7.2,-0.92) -- (7.4,-0.92);
  \draw[thick] (7.24,-1.02) -- (7.36,-1.02);
  \node[lbl] at (6.0,-0.45) {$\tC$};
  \node[lbl,text=black!45] at (7.0,-1.35) {legal: archived ($\gsink$)};
  % legal: reduced
  \node[box] (r) at (7.3,-2.3) {$\mathit{red}$};
  \draw[wire] (6.0,-2.3) -- (r.west);
  \draw[wire] (r.east) -- (8.4,-2.3);
  \node[lbl] at (6.0,-2.05) {$\tS$};
  \node[lbl] at (8.55,-2.05) {$\tA$};
  \node[lbl,text=black!45] at (7.3,-2.85) {legal: reduced};
\end{tikzpicture}
\caption{Negative pedagogy for the two laws. The red crosses mark shapes the
grammar simply does not generate: (a) there is no junction node for thin
wires, so ``$\tC \glto \tC\ten\tC$'' cannot be drawn (the knower rail
branches freely, by contrast); (b) a content wire cannot end in mid-air, and
it cannot be discarded — it must be structurally carried in the conclusion,
archived into $\gsink$ (a legal, terminal destination), or consumed by a
reducer. In proof nets such conditions are global correctness criteria; here
well-formedness is compositional, so the forbidden shapes are exactly the
undrawable ones (\S\ref{sec:laws}).}
\label{fig:laws}
\end{figure}


\begin{itemize}
\item \textbf{No cloning.} A thin wire has no junction. The shape ``one
$\tC$ wire in, two out'' requires a junction node, and the grammar has
none for thin wires. The unrestricted statement
\dcn{no_contraction_content} is a named \conjl{} in the
mechanization, proven in fragment strengths:
\begin{itemize}
\item the multiplicative core --- \dcn{no_contraction_core}
  \meas{machine-checked};
\item the generator-free fragment ---
  \dcn{no_contraction_nogen} \meas{machine-checked};
\item the pure-equality fragment ---
  \dcn{no_contraction_eq} \meas{machine-checked}.
\end{itemize}
The diagram grammar agrees with the fragments: within the core language
the shape is undrawable, and the full statement's remaining cases live
exactly where the calculus's own proof still lives.
\item \textbf{No vanishing.} A content wire has no bare end and no
discard terminator. It must be \emph{carried} (structurally present in the
conclusion), \emph{archived} (into $\gsink$, which is legal and terminal),
or \emph{reduced} (consumed by a consumer that exists in the signature).
The unrestricted statement \dcn{no_vanishing} is a named
\conjl{} with the strongest proven fragment ladder in the development:
\begin{itemize}
\item the equality fragment --- \dcn{no_vanishing_eq}
  \meas{machine-checked};
\item the clean fragment --- \dcn{no_vanishing_clean}
  \meas{machine-checked};
\item the parkfree fragment --- \dcn{no_vanishing_parkfree}
  \meas{machine-checked};
\item the conditional fixfree and fix-layer forms ---
  \dcn{no_vanishing_fixfree_conditional} and
  \dcn{no_vanishing_fix_conditional} \meas{machine-checked},
  each under its stated hypotheses.
\end{itemize}
\end{itemize}

Two scoping remarks, both easy to get wrong on first reading.

First: the laws are \emph{sorted}. The knower rail branches freely --- the
witness \dcn{knower_contractible} derives
$[\tK] \vdash \tK \ten \tK$ \meas{machine-checked} --- and that is not an
exception to no-cloning but its point: exactly one sort carries a
duplicable structure, and it is the sort that licenses the generator. What
the linear sorts deny is the \emph{uniform} comonoid; the companion paper
tracks the precise categorical reading
(\cite{dcarecompanion}, after \cite{abramsky2004}).

Second: the archive is not a loophole. The legal pattern
$\tC \glto \gsink$ is derivable --- the witness
\dcn{archival_witness} \meas{machine-checked} --- and it is the
calculus's most loaded legal shape: the count certifies the token was
\emph{spent} (the counting instance \dcn{no_weakening_content}
\meas{machine-checked}), while nothing certifies it was \emph{heard},
because nothing in the calculus reads a trace
(\dcn{past_cannot_teach} \meas{machine-checked}). The sink is
where a wire may legally end; that is precisely why the notation must
distinguish ending-in-$\gsink$ from ending-in-nothing, and why the second
is the forbidden shape, not the first.

\section{Two worked diagrams}\label{sec:worked}

\subsection{The epistemic cycle}\label{sec:worked:cycle}

% fig-cycle.tex — the epistemic cycle, drawn from Dcare.knower_sustains_stream.
\begin{figure}[t]
\centering
\begin{tikzpicture}
  % the knower source, OUTSIDE the envelope
  \node[circle,draw,very thick,fill=black!12,inner sep=0pt,minimum size=5mm]
       (kn) at (0,0) {$\cdot$};
  \node[lbl] at (0,-0.45) {$\gknower\;\tK$};
  % generator and cons inside the envelope
  \node[box] (gen) at (2.2,0) {$\mathit{gen}$};
  \node[box] (cons) at (6.2,0.35) {$\mathit{cons}$};
  % the thick K rail crosses the envelope boundary
  \draw[thickwire] (kn) -- (gen.west);
  % gen's outputs: thin C up, thick K down to a discard terminator
  \draw[wire] (gen.east) -- ++(0.45,0) -- (cons.north west);
  \node[lbl] at (4.1,0.75) {$\tC$ (head)};
  \draw[thickwire] (gen.east) -- ++(0.45,0) -- ++(0,-0.85)
        -- (3.6,-0.85);
  \node[discard] (disc) at (3.8,-0.85) {};
  \node[lbl] at (4.25,-0.85) {$\tK$ discarded (rail)};
  % cons output S
  \draw[wire] (cons.east) -- (9.0,0.35);
  \node[lbl] at (8.3,0.6) {$\tS$};
  % the feedback: from the output wire, down and around, through ONE bead,
  % into cons's tail port
  \node[bead] (fb) at (4.6,-2.1) {};
  \draw[wire] (7.4,0.35) -- (7.4,-1.3) -- (fb) -- (3.4,-2.1)
        -- (3.4,-1.0) -- (cons.south west);
  \node[lbl] at (6.0,-2.4) {$\gfeed$: the loop is a binding, and it passes
        through exactly one $\guard$};
  \node[lbl] at (5.55,-1.75) {$\guard\tS$ (tail)};
  % the envelope (drawn on the background so it sits under nothing)
  \begin{pgfonlayer}{background}
    \node[env,fit=(gen)(cons)(disc)(fb),inner sep=9pt] (env) {};
  \end{pgfonlayer}
  \node[lbl] at (5.1,1.35) {feedback envelope ($\mathit{fix}$): only the
        thick rail crosses it};
  % the reducer, OUTSIDE, consuming the stream
  \node[box] (red) at (10.1,0.35) {$\mathit{red}$};
  \draw[wire] (9.0,0.35) -- (red.west);
  \draw[wire] (red.east) -- (11.1,0.35);
  \node[lbl] at (11.25,0.6) {$\tA$};
\end{tikzpicture}
\caption{The epistemic cycle, drawn from the machine-checked witness of
\texttt{Dcare.knower\_sustains\_stream}: the term
$\mathit{fix}\,(\lambda \guard\tS.\;\mathit{lete}\;(\mathit{pair}(s,k))\;
\mathit{cons}(\mathit{extract}\ k,\ s))$ at context $[\tK]$. Each lap, the
generator fires on the knower and produces fresh content $\tC$ (its knower
half returns to the rail and is discarded — legal for a rail, illegal for a
wire); $\mathit{cons}$ assembles head and delayed tail; the output streams
out to the reducer. The feedback loop carries exactly one guard bead: one
lap $=$ one interval. Only the thick rail crosses the envelope — the
diagrammatic form of the side condition that makes self-application cost a
context.}
\label{fig:cycle}
\end{figure}

Figure~\ref{fig:cycle} draws the calculus's central positive result: a
live knower sustains an infinite guarded stream. In the Guarded Syntax
the cycle is one line --- the schematic form of the whole epistemic
loop:
\[
\gknower \;:\; \tS \glto \bigl(\guard\tA \gfeed \gknower\bigr)
\qquad\text{read: from the stream, guarded knowledge feeds back to the
knower, one $\guard$ per lap;}
\]
and the machine-checked witness
\dcn{knower_sustains_stream} \meas{machine-checked} exhibits its
realization at context $[\tK]$:
\[
s \gfeed x.\;\texttt{let}\;(k,c) = \mathit{gen}\ \gknower\ \texttt{in}\
  \mathit{cons}(c,\ x) \;:\; \tS
\]
where the feed's binder $x$ arrives already delayed (the binding
\emph{is} the $\guard$: one lap of the feed consumes exactly one
interval), and $\texttt{let}$ is the calculus's projection-free tensor
elimination. In the judgment's own spelling this is
$\mathit{fix}\bigl(\lambda \guard\tS.\;
\mathit{lete}\,(\mathit{pair}(s,k))\;
\mathit{cons}(\mathit{extract}\ k,\ s)\bigr) : \tS$, where
$\mathit{extract}$ fires the generator on $k$ and keeps the content
half (the knower half returns to the rail and is discarded --- legal for a
rail). Read the diagram one lap at a time. The thick rail crosses the
envelope --- the only wire that may --- and the generator fires on it,
producing fresh content now and the knower again. The content becomes the
head; the previous stream, delayed one interval, becomes the tail; the new
stream leaves the envelope for the reducer. The loop carries exactly one
bead: one lap is one interval. Generate, reduce, feed back --- the cycle
the calculus was built to type, as one closed picture whose every visual
feature corresponds to a clause of the rule it draws.

The figure also shows the two laws of \S\ref{sec:core:feedback} at work:
the envelope is permeable to the rail because the witness's context
$[\tK]$ satisfies the all-non-linear side condition
(\dcn{knower_ctx_nonlinear} \meas{machine-checked}), and the
loop is legal because it passes through the bead. Remove the bead from
the drawing and what remains is not a fragile diagram but an untyped one:
the grammar has no envelope for it.

\subsection{The cut, and the mutual promise}\label{sec:worked:cut}

% fig-cut.tex — the cut as splicing; the mutual promise as an unreduced composite.
\begin{figure}[t]
\centering
\begin{tikzpicture}
  % ---- panel (a): the splice ----
  \node[lbl,align=left] at (1.9,1.5) {(a) the cut: a pending\\ linear wire is spliced};
  % the promise box: context in, hole on the input fringe, conclusion out
  \node[box,minimum width=16mm,minimum height=14mm] (big) at (1.9,0.3) {};
  \node[font=\small] at (1.9,0.3) {$h$};
  \draw[wire] (0.0,0.75) -- (big.160);
  \node[lbl] at (-0.1,0.98) {$\Gamma_L$};
  \draw[wire] (0.0,-0.15) -- (big.200);
  \node[lbl] at (-0.1,-0.38) {$\Gamma_R$};
  \draw[wire] (big.east) -- (3.8,0.3);
  \node[lbl] at (3.95,0.55) {$\sigma$};
  % the hole: a dashed circle on the input fringe between the two blocks
  \node[draw,dashed,circle,inner sep=0pt,minimum size=3.4mm]
       (hole) at (0.55,0.3) {};
  \node[lbl] at (0.55,0.02) {$a$ (the hole)};
  % the argument box producing a
  \node[box] (arg) at (1.9,-1.7) {$\mathit{arg}$};
  \draw[wire] (arg.east) -- ++(1.5,0) -- ++(0,1.1) -- (3.8,-0.6);
  \node[lbl] at (4.25,-0.85) {};
  \node[lbl] at (3.6,-1.45) {$a$};
  % the splice arrow
  \node[font=\Large] at (5.3,0.0) {$\Rightarrow$};
  % after: one box, Delta inserted AT THE HOLE'S POSITION (ordered)
  \node[box,minimum width=16mm,minimum height=16mm] (fused) at (7.6,0.0) {};
  \node[font=\small] at (7.6,0.35) {$h'$};
  \node[font=\scriptsize\itshape] at (7.6,-0.35) {$\mathit{arg}$ spliced};
  \draw[wire] (5.9,0.55) -- (fused.160);
  \node[lbl] at (5.75,0.78) {$\Gamma_L$};
  \draw[wire] (5.9,-0.35) -- (fused.200);
  \node[lbl] at (5.75,-0.58) {$\Gamma_R$};
  \draw[wire] (5.9,0.1) -- (fused.185);
  \node[lbl] at (5.8,-0.05) {$\Delta$};
  \draw[wire] (fused.east) -- (9.5,0.0);
  \node[lbl] at (9.65,0.25) {$\sigma$};
  \node[lbl,text=black!45] at (7.6,-1.35)
        {$\Gamma_L \,{+}\, \Delta \,{+}\, \Gamma_R$, in that order};
  % ---- panel (b): the mutual promise ----
  \node[lbl,align=left] at (1.9,-3.3)
        {(b) the mutual promise:\\ two boxes, one ungueded cycle};
  \node[box] (l1) at (1.2,-4.3) {$\lambda_1$};
  \node[box] (l2) at (3.4,-4.3) {$\lambda_2$};
  \draw[wire] (l1.east) to[out=25,in=155] (l2.west);
  \draw[wire] (l2.east) to[out=-90,in=0] (4.4,-5.2) -- (0.2,-5.2)
        to[out=180,in=-90] (l1.west);
  \node[lbl,text=red!60!black] at (2.3,-5.55)
        {no $\guard$ anywhere on this cycle};
  \node[font=\Large] at (5.3,-4.6) {$\Rightarrow$};
  \node[box,minimum width=14mm] (nf) at (7.2,-4.6) {$\mathit{normal}$};
  \node[lbl,text=black!45] at (7.2,-5.3) {a ledger with no pending promises};
  \node[lbl] at (9.0,-4.6) {};
\end{tikzpicture}
\caption{(a) The cut. A derivation $h$ at an ordered context
$\Gamma_L \,{+}\, a \,{+}\, \Gamma_R$ leaves a hole of type $a$ on its input
fringe; the derivation of the argument plugs that hole, and the $\Delta$
block takes the hole's position — order is preserved, nothing is commuted.
This is the diagrammatic reading of \texttt{Dcare.linear\_cut}
\meas{machine-checked at the term level}, with the size law of
\texttt{Dcare.linear\_beta\_splice}: the spliced composite is exactly three
nodes smaller than the redex it resolves.
(b) The mutual promise — two abstractions, each waiting on the other's
output. It is \emph{not} a legal feedback (there is no guard bead on the
cycle, and thin wires cannot cross an envelope); it is an unreduced
composite, and the cut dissolves it.}
\label{fig:cut}
\end{figure}

The cut is the notation's substitution principle. A derivation $h$ at an
ordered context $\Gamma_L + a + \Gamma_R$ leaves, diagrammatically, a
\emph{hole} of type $a$ on its input fringe; a derivation of the argument
at $\Delta$ plugs that hole, and the $\Delta$ block takes the hole's
position in the order (Fig.~\ref{fig:cut}a). Nothing commutes: the splice
is the ordered-calculus cut, and the diagram shows the order preservation
as geometry --- the $\Delta$ wires enter where the hole was. At the term
level this is the machine-checked substitution principle
\dcn{linear_cut} \meas{machine-checked} --- a linear hypothesis
spent exactly once is replaced by its derivation, with full size and
counter accounting --- and its redex form
\dcn{linear_beta_splice} \meas{machine-checked}: the spliced
composite is exactly three nodes smaller than the redex it resolves.

The cut also resolves a shape that looks like feedback but is not
(Fig.~\ref{fig:cut}b): two abstractions, each waiting on the other's
output --- the \emph{mutual promise}. The cycle they form carries no guard
bead, and thin wires cannot cross an envelope, so the picture is not a
legal feedback; it is an unreduced composite. The cut dissolves it in both
directions at once, and what remains has no pending wires. A ledger with
no pending promises: that is the normal form of a mutual promise in this
calculus, and the diagram shows why it is a normal form and not a
deadlock --- there was never a loop, only two holes facing each other.

\section{The correspondence, stated precisely}\label{sec:correspondence}

This section states the diagram--derivation correspondence. All five
statements are \emph{claims}: the \texttt{Dcare} development proves the
term-level theorems cited alongside them, but the translation itself is
not mechanized, and no claim below should be read as machine-checked.
Table~\ref{tab:claims} summarizes.

\begin{claim}[soundness]\label{cl:sound}\claimS
Every well-formed diagram with input fringe $\Gamma$ and output $\tau$
denotes a term $t$ with a derivation $\mathit{Typed}\;\Gamma\;t\;\tau$.
\end{claim}
\begin{proofsketch}
Induction on diagram construction. Wires to ports give the variable and
constant rules; juxtaposition gives ordered context splitting (the split
equation records which side each box occupied); boxes give the signature
constants; a bead on a wire gives the delay; an envelope whose boundary is
crossed only by rails gives the guarded fixpoint, with the guard law
supplying the delayed binder and the permeability law supplying the
all-non-linear side condition; a swap-box unfolds to its defining
$\mathit{pair}$--$\mathit{lete}$ composite. The induction mirrors the rule
schema of the judgment case by case.
\end{proofsketch}

\begin{claim}[completeness]\label{cl:comp}\claimS
Every derivation draws: for each derivation of $\mathit{Typed}\;\Gamma\;t\;\tau$
there is a well-formed diagram with fringe $\Gamma$ and output $\tau$.
\end{claim}
\begin{proofsketch}
Induction on the derivation. The interesting case is the one that
motivates the swap-box: the tensor introduction binds its components in
one order and the elimination in the other, so a derivation that
re-orders two adjacent hypotheses draws as a $\mathit{pair}$ box feeding a
$\mathit{lete}$ box --- the definable exchange of
Fig.~\ref{fig:swap}, already exhibited inside the machine-checked witness
\dcn{knower_sustains_stream} \meas{machine-checked}. All other
rules are single boxes or wire compositions.
\end{proofsketch}

\begin{claim}[the cut]\label{cl:cut}\claimS
Diagrammatic splicing --- plugging a hole of type $a$ with a box
realizing $a$ --- corresponds to the ordered cut. Restricted to linear
splices (the hole's hypothesis spent exactly once) it is the term-level
\dcn{linear_cut} \meas{machine-checked}.
\end{claim}
\begin{proofsketch}
The diagram operation composes two boxes along one wire; at the term
level the machine-checked substitution principle performs exactly this
composition with size and counter preservation, and its redex form gives
the three-node size law drawn in Fig.~\ref{fig:cut}. The diagram-level
statement --- that graph rewriting along the splice is confluent and
terminating on cut-free forms --- is future work; the claim here is only
the correspondence of the one step.
\end{proofsketch}

\begin{claim}[guarded feedback]\label{cl:feed}\claimS
An envelope is well-formed if and only if every loop inside it passes
through at least one guard and its boundary is crossed only by rails. The
term-level reading of the two conditions is the guarded fixpoint rule:
the one-bead loop is the delayed binder (one-tick unfolding,
\dcn{fix_unfold_closed} \meas{machine-checked}) and the
permeability condition is the all-non-linear context side condition
(the typed counterexample \dcn{selfapp_price}
\meas{machine-checked} shows the condition is not removable).
\end{claim}
\begin{proofsketch}
Both directions are inductions on envelope structure. The guard law
prevents zero-delay cycles, which is the diagrammatic content of
guardedness as in synchronous dataflow causality \cite{lustre1991}; the
permeability law is a restatement of the fixpoint side condition, which
the counterexample shows cannot be dropped. What is \emph{not} claimed: a
denotational semantics of envelopes as a guarded trace
\cite{jsv1996,birkedal2012} --- that construction is the natural next
step and is open.
\end{proofsketch}

\begin{claim}[equivalence]\label{cl:eq}\claimS
Two diagrams denote the same term iff one is obtained from the other by
planar isotopy and the local moves that mirror the calculus's
computational rules (bead--box commutations, swap-box triangle laws,
splice). The swap-box's coherence laws hold as derived moves, not as
structure.
\end{claim}
\begin{proofsketch}
The forward direction is a case analysis on the derivation equality; the
swap-box laws reduce to the $\mathit{pair}$--$\mathit{lete}$ composite
that defines the box. The interesting part --- that no \emph{additional}
identifications hold, i.e.\ that crossing-free isotopy plus the moves is
complete for the ordered judgment --- is the honest gap; it is the
usual coherence question \cite{joyal1991,selinger2011} transplanted to
this setting, and we state it as open.
\end{proofsketch}

\begin{table}[t]
\centering
\begin{tabular}{llp{62mm}}
\textbf{claim} & \textbf{status} & \textbf{term-level anchor (machine-checked)}\\
\hline
\ref{cl:sound} soundness & claim, sketched & the rule schema of \texttt{Typed}\\
\ref{cl:comp} completeness & claim, sketched & \dcn{knower_sustains_stream} (the swap case)\\
\ref{cl:cut} cut & claim, one step & \dcn{linear_cut}, \dcn{linear_beta_splice}\\
\ref{cl:feed} guarded feedback & claim, sketched & \dcn{fix_unfold_closed}, \dcn{selfapp_price}\\
\ref{cl:eq} equivalence & claim, open & none (coherence question)\\
\end{tabular}
\caption{The correspondence at a glance. Every anchor names a proven
declaration in the \texttt{Dcare} development; the diagram-level
statements themselves are claims.}
\label{tab:claims}
\end{table}

\section{Related work}\label{sec:related}

\paragraph{Proof nets.} Girard's proof nets \cite{girard1987} and the
Danos--Regnier correctness criteria \cite{danos1989} give
multiplicative linear logic its geometry. The Guarded Syntax differs in
three deliberate ways: its wires are directed (dataflow, not proof
search); its well-formedness is compositional rather than a global graph
criterion; and delay and feedback are primitive rather than encoded. A
net-style semantics for the Guarded Syntax --- boxes as links, beads as
unit-delay links --- is a natural refinement of Claim~\ref{cl:eq}.

\paragraph{String diagrams for monoidal categories.} Joyal and Street's
geometry of tensor calculus \cite{joyal1991} and Selinger's survey
\cite{selinger2011} fix the vocabulary: boxes, wires, isotopy. The
Guarded Syntax takes the \emph{planar} fragment --- the non-symmetric
case, matching ordered logic \cite{polakow2017} --- and adds a directed
flow, a wire decorator for delay, and a guarded envelope. Of the standard
machinery, the envelope is closest to the trace of a traced monoidal
category \cite{jsv1996}, with the guard law playing the role of the
axiom that makes the trace exist.

\paragraph{Interaction nets.} Lafont's interaction combinators
\cite{lafont1997} make cut elimination itself the computation, as local
graph rewriting. The splice of Fig.~\ref{fig:cut} is exactly an
interaction step seen from the ordered-linear side; an interaction-net
presentation of the Care signature, with the guard as a scheduler, is a
plausible implementation substrate for the notation.

\paragraph{Signal-flow graphs.} The prop semantics of signal-flow graphs
\cite{bsz2014} and synchronous dataflow languages \cite{lustre1991} are
the engineering kin: directed wires, boxes, and feedback loops whose
causality is guaranteed by delays. Their wires fan out freely (a
Cartesian setting); the Guarded Syntax is the same picture under a
linearity discipline, where fan-out is the one thing a wire cannot do.

\paragraph{Guarded recursion.} The delay bead is the diagrammatic form of
the later modality: Nakano's modality for recursion \cite{nakano2000}
and the topos-of-trees line \cite{birkedal2012} supply the semantics of
$\guard$ and of guarded fixpoints. The permeability law is the
diagrammatic shadow of the side conditions those systems place on
recursive definitions.

\paragraph{The Care Calculus.} The calculus, its mechanization, and its
programmatic context are the companion paper's subject
\cite{dcarecompanion}; the present paper depends on it for every
term-level theorem cited, and contributes only the notation and the
correspondence statements.

\section{Interpretive remarks}\label{sec:remarks}

The notation was built for proof-reading, but three of its features have
readings beyond proof theory. Each is labeled; none is a result of this
paper or of the mechanization.

\paragraph{Discreteness.} The guard bead gives every feedback loop a
clock: one lap, one interval, and no lap without one. The notation
therefore has \emph{discrete time in its vocabulary} --- but a vocabulary
is not a physics; whether any physical process satisfies the grammar's
clock is a question the notation equips one to ask and does not answer.
It is a reading, not a result.

\paragraph{Conservation at the boundary.} The two ungeneratable shapes
are a conservation pair: what cannot be copied and what cannot be
deleted. The same pair organizes quantum information theory's
no-cloning and no-deleting results
\cite{wootters1982,pati2000}, and the categorical semantics there
\cite{abramsky2004} explains why: duplicability and deletability are
comonoid structure, and the resource sorts deny it. The companion paper
develops the identification and its precise scope (a naturality
condition, without which the reading overclaims)
\cite{dcarecompanion}. The Guarded Syntax contributes only the
observation that in two dimensions the pair becomes visual: a wire that
cannot split and a wire that cannot stop. It is a reading, not a result.

\paragraph{Rotation.} The cycle of Fig.~\ref{fig:cycle} reads the same
from any position on the loop: each lap generates, reduces, and feeds
back, and no point on the loop is privileged. This is an observation
about the diagram's symmetry group, and it is tempting to say more about
systems whose diagrams look like this; the notation itself does not
license more. It is a reading, not a result.

\section{Conclusion}\label{sec:conclusion}

The Guarded Syntax is a planar, directed, two-gauge string-diagram
notation in which the resource discipline of an ordered linear calculus
--- what is consumed, what is delayed, what loops, what is archived ---
is visible on the page. Its grammar is compositional, so the calculus's
two linearity laws are ungeneratable shapes rather than audited
conditions; its feedback is guarded, so every loop has a clock; its
cut is a splice, so circular obligation dissolves into normal form. The
correspondence with the mechanized judgment is stated as five claims,
anchored at the term level to machine-checked theorems and honest about
what remains open: the coherence question, the guarded-trace semantics,
and a mechanized diagram checker, which is the natural next unit of
work.

\subsection*{Acknowledgements} The notation was developed against a
mechanization whose every rule was audited by proof pressure; the
diagrams in this paper are drawn from its witnesses.

\bibliographystyle{eptcs}
\bibliography{refs}

\end{document}
